\chapter{System Implementation and Methodology}

\section{Technology Stack}
The backend is implemented in Python using FastAPI, Pydantic, SQLAlchemy, pandas, NumPy, scikit-learn, XGBoost, SHAP, LIME, and ReportLab. The frontend is a React application built with Vite and styled with Tailwind CSS. Axios is used for API calls. The main application is organized around one loan-prediction router and a registry that loads the home-loan, personal-loan, and legacy model pipelines. The deployed product paths are the home-loan and personal-loan pipelines; the legacy pipeline remains for backwards compatibility.

\section{Development Environment}
The repository README documents Python 3.10 or later, Node.js 18 or later, and separate backend and frontend development servers. The API is typically served by Uvicorn on port 8000 and Vite serves the client during development. Dependencies are described in \texttt{requirements.txt} and \texttt{frontend/package.json}. The Python requirements use lower version bounds rather than a fully frozen environment, so exact reproducibility also depends on the installed package versions and platform.

\section{Dataset Description}
Two product datasets are present: \path{datasets/homeloan_train.csv} and \path{datasets/personalloan_train.csv}. Repository scripts generate each dataset with 3,000 records using NumPy random seed 42. The home-loan table has 18 columns and the personal-loan table has 17. Both contain a loan ID, applicant attributes, financial inputs, credit information, and a target label named \texttt{Loan\_Status}. The home-loan dataset has 1,529 approved and 1,471 rejected records. The personal-loan dataset has 1,563 approved and 1,437 rejected records. The checked CSVs have no missing cells and no duplicate full rows.

These are synthetic data generated by the repository. Their label-generation scripts construct the outcome from hand-specified financial and credit rules and include a gender-dependent adjustment in a borderline region to create a fairness-analysis example. The high test scores reported later are consequently measurements on a designed simulation, not evidence of performance on a bank's historical or future applicant population. There is no externally sourced lending dataset or data dictionary in the repository.

\section{Data Preprocessing}
The training scripts read a product CSV, compute derived features, separate \texttt{Loan\_Status} from the predictors, and drop \texttt{Loan\_ID}. The target is mapped to 1 for Approved and 0 for Rejected. An 80/20 stratified train-test split is created with random state 42. The preprocessor uses a column transformer: categorical fields pass through a one-hot encoder with the first category dropped and unknown values ignored; numerical fields pass through standard scaling. The transformer is fit on the training set and then applied to the test set.

For online requests, Pydantic schemas normalize optional numeric values and the model registry fills product-specific defaults. This normalization is operationally convenient but means omitted values can be replaced with defaults; users should provide valid application data and review the displayed submitted form. Batch processing applies row-level normalization and records rows that cannot be used as flawed records with explanations of the validation issue.

\section{Exploratory Data Analysis}
No reproducible exploratory-data-analysis notebook or EDA report is included in the current source tree. The dataset-generation scripts specify distributions for demographics, income, loan terms, property characteristics, CIBIL scores, utilization, and missed payments. This section therefore describes those generator rules and the observed row/label counts rather than claiming a separate EDA study. A manually produced chart of class counts may be added at \texttt{figures/graphs/dataset\_class\_distribution.png} as shown in Graph~\ref{graph:dataset-classes}.

\reportgraph{figures/graphs/dataset_class_distribution.png}{0.78}{3.2}{Generated dataset class counts}{Approved and rejected counts in the two generated product datasets; use the repository CSV counts reported in this chapter.}{graph:dataset-classes}

\section{Feature Engineering}
The two products use distinct financial calculations. For a loan principal $P$, annual interest rate $i$, and term $n$ months, let $r=i/12$. The amortized monthly instalment used by the product pipelines is
\begin{equation}
  \mathrm{EMI}=P\frac{r(1+r)^n}{(1+r)^n-1}.
  \label{eq:emi}
\end{equation}
The home-loan generator and preprocessing function use an illustrative annual rate of 8.5\%; the personal-loan path uses 13.5\%. These fixed rates are implementation assumptions for engineered features, not live lender quotations.

For a home loan, monthly household income is the sum of applicant and co-applicant income. For a personal loan, the FOIR denominator is applicant income. In both product-specific pipelines the ratio is
\begin{equation}
  \mathrm{FOIR}=\frac{\mathrm{Existing\ EMI}+\mathrm{Proposed\ EMI}}{\mathrm{Income\ used\ for\ product}}.
  \label{eq:foir}
\end{equation}
The home-loan pipeline also calculates
\begin{equation}
  \mathrm{LTV\ Ratio}=100\times\frac{\mathrm{Loan\ Amount}}{\mathrm{Property\ Value}}.
  \label{eq:ltv}
\end{equation}
Personal-loan features include loan purpose and credit-card utilization instead of property collateral features. Home-loan raw fields include the property value and area; personal-loan raw fields include purpose and utilization. Engineered EMI/FOIR and product-specific columns are included in each model's numerical feature set.

\section{Machine Learning Model Development}
Separate XGBoost classifiers are trained in \path{scripts/train_homeloan.py} and \path{scripts/train_personalloan.py}. Both scripts use 120 estimators, maximum tree depth 4, learning rate 0.08, row subsampling 0.85, column subsampling 0.85, random state 42, and log-loss as the evaluation metric. Each script saves the classifier, fitted preprocessor, transformed feature names, a 100-row transformed background sample, and metrics under the product directory. The runtime registry loads those artifacts into a \texttt{ProductPipeline} with the product feature function and explainers.

\subsection*{Algorithm 1: Product-specific model training}
\begin{lstlisting}
Input: Product dataset D, product feature function g, seed 42
Output: Fitted preprocessor T, XGBoost classifier M, test metrics
1. D' <- g(D); remove Loan_ID from predictors
2. Encode Loan_Status: Approved=1, Rejected=0
3. Split predictors and labels with stratification (80% / 20%)
4. Fit one-hot encoding and standard scaling on training predictors
5. Transform training and test predictors with the fitted transformer
6. Fit XGBoost(n_estimators=120, max_depth=4, learning_rate=0.08,
   subsample=0.85, colsample_bytree=0.85, random_state=42)
7. Predict labels and approval probabilities on the held-out test split
8. Save model, transformer, feature names, background sample, and metrics
\end{lstlisting}

\section{Model Training and Evaluation}
The training scripts calculate accuracy, precision, recall, F1-score, and ROC-AUC on the held-out split. Product-specific metrics are saved in \texttt{models/homeloan/metrics.json} and \texttt{models/personalloan/metrics.json}. No product-specific confusion matrix is stored in those files. The root \texttt{models/model\_metadata.json} describes an older legacy model and must not be substituted for the current product metrics. Chapter 7 reports the product-specific artifacts and keeps the legacy results separate.

\section{SHAP Implementation}
The model registry initializes a \texttt{SHAPExplainer} backed by \texttt{shap.TreeExplainer} for each XGBoost classifier. For each application the transformed feature vector is passed to the explainer. The implementation extracts the contribution vector, maps the sign to positive, negative, or neutral direction, and returns one feature record for each transformed model feature. Local contribution values are in the SHAP output space used by the explainer; they are not percentages of approval probability. Displaying raw values for one-hot and engineered features requires care because the feature name may not match a raw applicant field.

\section{LIME Implementation}
The \texttt{LIMEExplainer} wraps \texttt{lime.lime\_tabular.LimeTabularExplainer} in classification mode. The explainer receives the transformed training background and model feature names. When explaining an input, it calls the classifier's \texttt{predict\_proba} function and requests contributions for the approved class. The route sets a deterministic seed derived from the input JSON using SHA-256 before running LIME. This supports repeatable sampling for the same input in a fixed software environment; it does not imply that all library versions or hardware produce bitwise identical results.

\section{Explanation Validation Framework}
The route applies four checks after generating the local explanation. The formulas below reproduce the code-level definitions and are not presented as universal standards.

\subsection{Faithfulness}
The implementation uses the SHAP explanation for the current model prediction. For an approved prediction, it selects positive SHAP contributors above 0.08, falling back to up to five largest positive contributions. For a rejected prediction, it selects negative contributors below $-0.08$, falling back to up to five largest negative contributions. The selected transformed coordinates are replaced with zero. Let $p$ be the original approval probability, $p_m$ the probability after masking, and $p_0$ the probability for the zero-vector baseline. Define
\begin{equation}
  \Delta=|p-p_m|,\qquad q=\frac{\Delta}{\max(0.08,|p-p_0|)}.
\end{equation}
The implemented score is
\begin{equation}
 F=\begin{cases}
 \max(0.02,2q), & q\leq 0.05,\\
 \min(0.95,0.74+0.17(1-e^{-q})), & q>0.05.
 \end{cases}
 \label{eq:faithfulness}
\end{equation}
If no salient features are selected, the function returns 0.10. The piecewise mapping is a project calibration heuristic. Its numerical scale is not a conventional probability of explanation correctness and should not be compared directly with unrelated faithfulness measures.

\subsection{Stability}
The route perturbs the transformed input with a Gaussian random vector of standard deviation 0.01 using a seed derived from the applicant input. The original and perturbed SHAP contribution vectors are compared with Spearman rank correlation $\rho$. The returned score is
\begin{equation}
  S=\frac{\rho+1}{2},
\end{equation}
which maps the correlation range from $[-1,1]$ to $[0,1]$. If the correlation is undefined because contributions are constant, the implementation returns zero. This measures rank similarity under the specific perturbation procedure; it is not a guarantee of invariance under all plausible applicant changes.

\subsection{Consistency}
Let $A_3$ and $B_3$ be the sets of three features with largest absolute contributions under SHAP and LIME. The consistency score is their Jaccard similarity:
\begin{equation}
  C=\frac{|A_3\cap B_3|}{|A_3\cup B_3|}.
  \label{eq:consistency}
\end{equation}
The code compares feature names only; it does not compare contribution magnitudes or directions.

\subsection{Individual Counterfactual Fairness Score}
For a normalized applicant profile $x$, the code obtains approval probabilities $p(x)$ and $p(x^{(g)})$, where $x^{(g)}$ has the gender field flipped and other fields held fixed. The score is
\begin{equation}
  R=\max\left(0,\min\left(1,1-|p(x)-p(x^{(g)})|\right)\right).
  \label{eq:counterfactual}
\end{equation}
This is an individual model-sensitivity check for one binary attribute. It is not an estimate of group approval-rate parity, equal opportunity, causal fairness, or legal compliance.

\section{Trust Score Calculation}
The Trust Score is the weighted sum below, with all component scores on the interval $[0,1]$:
\begin{equation}
  T=0.35F+0.25S+0.30C+0.10R.
  \label{eq:trust-score}
\end{equation}
The weights are stored in \texttt{backend/app/trust/trust\_score.py}. Faithfulness receives 35\%, stability 25\%, SHAP-LIME consistency 30\%, and the counterfactual gender-flip score 10\%. The result is a project-defined summary used for presentation. No external calibration study or validated decision threshold for this composite is included.

\section{Confidence Score}
Let $p$ be the model's probability for approval and $\hat y$ its predicted class. The response confidence value is
\begin{equation}
  Q=\begin{cases}p,&\hat y=1,\\1-p,&\hat y=0.\end{cases}
\end{equation}
It is the model probability assigned to the predicted class. The repository does not run a separate calibration method such as isotonic regression or Platt scaling for these product models. Therefore the value should be called a prediction-confidence representation, not a verified calibration estimate or an explanation-reliability measure.

\section{Backend Implementation}
The FastAPI router is mounted under \texttt{/api/v1/loan}. \texttt{POST /predict} performs a single evaluation. Batch endpoints accept CSV uploads or JSON rows. Other endpoints list, retrieve, and delete evaluation or batch history, return sample CSVs and presets, and generate a PDF report. Pydantic schemas define application fields, optional product fields, explanation structures, solvency data, and response objects.

\section{Frontend Implementation}
The React application uses a screen state to select a portal mode, loan product, application form, and result view. Bank-officer and applicant screens are different presentations over shared evaluation results. The dashboard renders prediction information, SHAP/LIME summaries, validation metrics, submitted form values, and a loan-specific what-if affordability calculator. The what-if calculator recomputes EMI and FOIR with fixed illustrative rates; it does not call the model on slider changes. It should therefore not be read as a simulated approval prediction.

\section{Database Implementation}
SQLAlchemy models \texttt{PredictionLog} and \texttt{BatchUpload} persist evaluation and batch information. Local SQLite is always initialized as a fallback; a PostgreSQL URL can be supplied through \texttt{DATABASE\_URL}. The request session attempts the primary engine and falls back to SQLite after a connection failure. Database values include JSON fields for raw inputs, full responses, explanations, and flawed batch data. This makes the project runnable without a remote database, although local SQLite and cloud PostgreSQL are operational alternatives rather than a replicated multi-database consistency design.

\section{API Communication}
The frontend sends JSON requests for single predictions and PDF creation, and multipart form data for CSV uploads. The response includes prediction, approval probability, predicted-class confidence representation, model version, SHAP and LIME feature arrays, Trust Score components, solvency summary, and dossier ID. The API's structured response makes each output separately available to the UI and stored history.

\section{Authentication and Security}
The system does not implement account creation, passwords, tokens, OAuth, or server-side role authorization. Applicant and bank-officer selections change client-side navigation and visible controls. The backend CORS middleware allows all origins and credentials in the current configuration. These security properties define the system's deployment limits when handling personal financial data. The project does not claim security certification or regulatory compliance.

\section{PDF Report Generation}
The reporting module uses ReportLab to create an A4 PDF with application reference, applicant and facility summary, decision rationale, solvency information, trust and validation values, and next steps. The PDF endpoint receives applicant data and the prediction response and returns a PDF response. The report is generated on demand; the code does not create a cryptographic document hash or a signed officer approval record.

\section{End-to-End System Workflow}
\subsection*{Algorithm 2: Online loan evaluation}
\begin{lstlisting}
Input: LoanApplicationRequest x
Output: prediction response y
1. Select product pipeline M from x.loan_type
2. Normalize product-specific values and construct one-row data frame
3. Derive EMI, FOIR, and (for home loans) LTV features
4. Transform the row with the product's fitted preprocessor
5. Obtain model class y_hat and approval probability p
6. Generate SHAP and LIME local explanations
7. Calculate faithfulness, stability, consistency, and gender-flip scores
8. Calculate weighted Trust Score and predicted-class probability value
9. Generate product-prefixed dossier ID
10. Persist record when database session is available; return response
\end{lstlisting}

This sequence documents the behavior in the API route. It also shows why an explanation-quality score should remain distinct from the prediction probability and the deterministic FOIR summary.
